\section{Objetivo de entrenamiento y rollout: del aprendizaje supervisado al control en lazo cerrado}

El rasgo de ingeniería más atractivo de Decision Transformer es que durante el entrenamiento usa una action prediction loss estándar. Se puede entrenar gracias a un objetivo supervisado conocido, pero el despliegue sigue siendo un rollout en lazo cerrado: el modelo emite una acción, el entorno produce un nuevo estado y una reward, y el sistema actualiza el objetivo restante y vuelve a invocar el modelo. Esta sección separa estrictamente el entrenamiento de la inferencia.

\subsection{Loss para acciones continuous y discrete}

El objetivo de entrenamiento depende de la forma estadística del action space, pero ambos casos comparten el mismo principio: tomar directamente la logged action como label supervisada, en lugar de construir primero un value target. El control continuo suele hacer regresión sobre un vector de números reales, mientras que el control discreto predice una distribución de probabilidad sobre las action categories. Al leer las fórmulas conviene revisar al mismo tiempo la action normalization, la mask de timesteps válidos y el modo de reduction, porque estos detalles de implementación cambian el peso que cada episode tiene en el gradiente. Para una continuous action puede usarse el error cuadrático medio:
\[
\mathcal{L}_{\mathrm{MSE}}(\theta)
=\frac{1}{T}\sum_{t=1}^{T}\lVert a_t-\widehat{a}_t\rVert_2^2.
\]
Para una discrete action puede usarse cross-entropy:
\[
\mathcal{L}_{\mathrm{CE}}(\theta)
=-\frac{1}{T}\sum_{t=1}^{T}\log p_{\theta}(a_t\mid \widehat{R}_{\le t},s_{\le t},a_{<t}).
\]
$a_t$ es la action real del conjunto de datos, $\widehat{a}_t$ o $p_{\theta}$ es la salida del modelo y $T$ es la longitud efectiva de la secuencia. Las posiciones de padding deben excluirse mediante una mask; de lo contrario, los episodes cortos quedan contaminados por tokens artificiales.

\lecturefigure{slide-11-loss-function.jpg}{Objetivo de entrenamiento: la continuous action usa MSE y la discrete action usa cross-entropy.}{Video oficial de Stanford Online, 23:40.}

\begin{knowledgebox}{Lectura de la figura: el verdadero argumento es «No dynamic programming»}
La conclusión en negritas al pie de la slide importa más que las dos fórmulas. El RL tradicional value-based suele usar un Bellman target, y ese target depende a su vez de otro valor estimado o de una slowly updated network; Decision Transformer ajusta directamente las action labels del registro, de modo que su interfaz de entrenamiento se parece más a supervised learning. El costo es que el modelo no realiza policy improvement de forma automática: la mejora tiene que venir del return conditioning, de la combinación de datos y de la generalization del modelo.
\end{knowledgebox}

\teachervoice{El docente enfatiza que el valor de MSE y de cross-entropy está en que son «estables y fáciles de regularizar», no en que sean más prodigiosas que una RL loss. Reformular un problema complejo como aprendizaje supervisado permite reutilizar infraestructura madura, pero también deja el éxito o el fracaso más en manos de la calidad de los datos y del diseño del condicionamiento.}

\subsection{Cómo se actualiza el return-to-go durante la inferencia}

Sea $R_{\mathrm{target}}$ el objetivo inicial. En el instante $t$, el modelo muestrea o emite la acción $a_t$ a partir del contexto actual; el entorno devuelve la reward $r_t$ y el nuevo estado $s_{t+1}$. En el caso sin descuento se actualiza:
\[
\widehat{R}_{t+1}=\widehat{R}_t-r_t.
\]
Si se usa un return con descuento, hay que mantener la misma definición que en el entrenamiento, en lugar de copiar sin más la resta. En cada paso también hay que agregar los nuevos tokens al contexto y conservar solo los $K$ timesteps más recientes.

\lecturefigure{slide-12-rollouts.jpg}{El rollout se parece a la autoregressive generation: se fija un retorno objetivo, se genera una acción, se agrega el nuevo estado y se actualiza el retorno restante.}{Video oficial de Stanford Online, 24:58.}

\begin{lstlisting}[caption={Rollout autoregressive con return-to-go sin descuento},label={lst:dt-rollout}]
target_return = desired_return
state = env.reset()
returns, states, actions, timesteps = [target_return], [state], [], [0]

while not done:
    action = model(returns, states, actions, timesteps)[-1]
    next_state, reward, done, info = env.step(action)

    actions.append(action)
    states.append(next_state)
    returns.append(returns[-1] - reward)
    timesteps.append(timesteps[-1] + 1)

    returns, states, actions, timesteps = keep_last_k_steps(
        returns, states, actions, timesteps, K
    )
\end{lstlisting}

\subsection{Un ejemplo numérico}

Supongamos que el retorno objetivo es 10. Si las rewards reales de los tres primeros pasos son $2,0,3$, el remaining return que recibe el modelo es, en orden, $10,8,8,5$. Si el segundo paso no da reward, el objetivo no disminuye por sí solo; si algún paso obtiene una reward mayor de lo esperado, el objetivo restante se acerca a cero más rápido. Este proceso hace que el modelo sepa en todo momento «cuánto falta», pero supone que en los datos de entrenamiento haya suficientes combinaciones de objetivos restantes y estados similares.

\begin{importantbox}{Los cuatro invariantes del rollout}
\begin{enumerate}
  \item El entrenamiento y la inferencia deben usar la misma definición de return y la misma reward scale;
  \item la action actual no puede filtrarse a la posición de entrada que la predice;
  \item lo que se actualiza es el objetivo restante después de la reward realmente obtenida, no la reward que predice el propio modelo;
  \item la context truncation debe alinearse por timesteps completos; no se pueden partir las ternas $(R,s,a)$.
\end{enumerate}
\end{importantbox}

\begin{warningbox}{La exposure bias sigue presente}
Durante el entrenamiento, el modelo ve las actions históricas reales del registro; durante el rollout, ve sus propias actions históricas. Un error pequeño cambia el state futuro y provoca una desviación respecto de la distribución de entrenamiento. Que el objetivo supervisado de Decision Transformer sea estable no significa que la closed-loop error accumulation esté resuelta.
\end{warningbox}

\subsection{Resumen de la sección}

En la fase de entrenamiento se ajusta directamente la action, lo que evita el Bellman backup explícito; en la fase de inferencia se interactúa con el entorno mediante un autoregressive loop y se actualizan de forma continua el remaining return y el context. El aprendizaje supervisado simplifica la optimización, pero no elimina el desplazamiento de distribución ni la acumulación de errores durante el rollout.
